\documentclass[11pt,a4paper]{article}
\usepackage{fontspec,amsmath,amssymb,amsthm,microtype}
\setmainfont{lmroman10-regular.otf}[BoldFont=lmroman10-bold.otf,ItalicFont=lmroman10-italic.otf,BoldItalicFont=lmroman10-bolditalic.otf]
\usepackage[margin=25mm]{geometry}
\usepackage[unicode,hidelinks]{hyperref}
\hypersetup{pdfauthor={LangevinLean Project}}
\setlength{\emergencystretch}{3em}
\setlength{\parskip}{0.4em}
\setlength{\parindent}{0pt}
\allowdisplaybreaks[2]
\raggedbottom
\title{\texorpdfstring{PIPGLA about the original center for every positive proximal scale}{PIPGLA about the original center for every positive proximal scale}}
\author{LangevinLean Project}
\date{}
\begin{document}\maketitle
Draft mathematical exposition; these notes may contain mistakes.


The equal-step PIPGLA transport theorem supplies a convergent reference estimate with finite coefficient. Pathwise sensitivity compares any positive proximal scale with that equal-step chain. Projection and the finite smoothing displacement then give the original-center RMS bound for the actual splitting algorithm, with the scale-mismatch cost shown explicitly.




Assume the same strongly convex smooth \(g_1\), convex bounded-below \(K\)-Lipschitz \(g_2\), \(N>0\), \(0<m\le L_1\), \(\lambda,\gamma>0\), and \(\gamma\le2/(m+L_1)\). Initialization is in \(L^2\) and the actual coordinate Gaussian innovations satisfy joint independence.


\section{\texorpdfstring{Convert the pathwise scale comparison into RMS}{Convert the pathwise scale comparison into RMS}}


Drive the \(\lambda\)- and \(\eta\)-chains by the same initial variable and innovations. Their state distance is at most \(K|\lambda-\eta|/(m\gamma)\). Both chains have finite second moments. The parameter projection is nonexpansive, so Minkowski gives

\[\operatorname{RMS}_\lambda(a)\le\operatorname{RMS}_\eta(a)+\frac{K|\lambda-\eta|}{m\gamma}.\]

The proof works with the actual \(L^2\) realizations and then identifies their pushforward integrals.




Here is the conversion from a pathwise comparison to the two displayed law integrals. On the supplied probability space let \(X\) and \(Y\) be the actual \(L^2\) PIPGLA iterates with proximal parameters \(\lambda\) and \(\eta\), the same initialization, and the same innovations. The pathwise parameter-sensitivity theorem gives

\[\|X(\omega)-Y(\omega)\|\le c:=\frac{K|\lambda-\eta|}{m\gamma}\quad\text{almost surely}.\]



The pathwise bound is uniform over the probability source, so its \(L^2\) norm is at most the same constant because the source mass is one. Projection cannot increase that norm. Adding the equal-reference-center error by Minkowski therefore gives exactly the displayed RMS comparison. Both laws are represented by actual square-integrable chain iterates, ensuring the square-root integrals are finite.



Earlier results used here: \cite[Section~3.1.1]{LangevinLean2026IPLA}; \cite[Section~3.1.6]{LangevinLean2026IPLA}.

\section{\texorpdfstring{Combine equal-\allowbreak{}step transport and original-\allowbreak{}center projection}{Combine equal-step transport and original-center projection}}


Set the comparison scale to \(\eta=\gamma\). The equal-step law theorem and the original-center projection inequality give

\[\begin{aligned}\operatorname{RMS}_\lambda(\theta_U)\le{}&\frac{K|\lambda-\gamma|}{m\gamma}+\sqrt{\frac{\gamma K^2}{m}}+\sqrt{\frac{d_\theta}{Nm}}\\&+\sqrt{(1-m\gamma)^nW_2(\rho_0,\Pi_{N,\gamma})^2+\frac\gamma N C_N}.\end{aligned}\]

Here \(C_N\) is the explicit finite coefficient defined in the transport estimate. It is independent of the time step. Every positive \(\lambda\) is covered, with the displayed parameter-comparison cost; a vanishing error as \(\gamma\to0\) additionally requires controlling that cost. This is the proved replacement rate, not the paper’s printed heat-flow constants.




The equal-step transport theorem controls the actual law at proximal scale \(\gamma\) (with the lifted smoothing scale prescribed in its statement). The projection theorem adds the sharp equilibrium RMS and the original-center displacement. To handle a different positive proximal parameter \(\lambda\), compare the two pointwise runs under the same innovations: their uniform difference is at most \(K|\lambda-\gamma|/(m\gamma)\). Integrating this deterministic bound yields the same RMS difference because the probability has mass one. The final result adds this scale-comparison term to the already proved equal-step error; the parameter corridor then supplies its explicit simpler upper bound.





Take the reference proximal scale equal to the physical forward step. Its squared transport bound gives a square root of the transient-plus-floor expression. Projection adds sharp equilibrium spread and original-center smoothing displacement. Finally, scale sensitivity adds \(K|\lambda-\gamma|/(m\gamma)\). Covering all positive \(\lambda\) does not imply a uniform vanishing small-step bound unless this extra mismatch term is controlled as the step decreases.



Earlier results used here: \cite[Section~3.17.4]{LangevinLean2026ULASC}; \cite[Section~3.1.1]{LangevinLean2026IPLA}; \cite[Section~3.1.6]{LangevinLean2026IPLA}; \cite[Section~3.3.1]{LangevinLean2026IPLA}; \cite[Section~3.3.2]{LangevinLean2026IPLA}.

\begin{thebibliography}{ULA-LSI}
\bibitem[ULA-SC]{LangevinLean2026ULASC}
LangevinLean Project. \emph{Unadjusted Langevin under Strong Convexity}. LL-01, 2026. Draft report, content revision 8d6335dce333. \url{https://odakyildiz.github.io/LangevinLean/docs/pdf/UnadjustedLangevin/FullProof.pdf}.
\bibitem[IPLA-SC]{LangevinLean2026IPLA}
LangevinLean Project. \emph{Interacting Particle Langevin under Strong Convexity}. LL-03, 2026. Draft report, content revision c72e54460954. \url{https://odakyildiz.github.io/LangevinLean/docs/pdf/IPLA/FullProof.pdf}.
\end{thebibliography}

\end{document}
